\section{Introduction}

% Competition is a strong deterant for companies, research labs, and governmental insitutions to exchange
% resources. For example, pharaceutical companies may be hesitant to share data on different classes of 
% drugs with their competitors for fear of losing intellectual property. Similarly, research labs 
% may avoid sharing their equipment or findings,
% fearing that other labs would reap the rewards of publishing first. 
% However, it is often socially desirable for these insitutions to collaborate. In the aforementioned 
% examples, collaboration could expidite
% the development of novel pharaceuticals and scientific discoveries, both of which are beneficial
% to society.
% Moreover, while competition may deter the exchange of resources, it can be still be 
% mutually beneficial in non-zero sum settings.
% An important question is thus how to incentivize two competitors to share resources?

Competition often deters companies, research labs, and government institutions from sharing resources. 
For example, pharmaceutical companies may hesitate to share data on drug classes with competitors 
for feat of losing intellectual property. Likewise, research labs might avoid sharing equipment or findings to prevent 
others from publishing first. Despite this, collaboration is often socially beneficial. In the aforementioned cases, 
it could accelerate the development of new drugs and scientific discoveries. 
Moreover, while competition can hinder cooperation, resource sharing may still be mutually advantageous 
in nonzero sum settings. 

An obstacle to resource sharing between two competitors is that each may be uncertain of the value of the 
other's resources until after using them. For example, the utility of a dataset might only become 
clear once it's used to train a machine learning model. Similarly, the value of access to compute or equipment 
may depend on whether the resulting experiments or simulations produce novel findings. 
% This work explores how to incentivize 
The central question this work answers is how to incentivize
two competing agents to share resources when mutual benefit is likely, but their valuations for each 
others' resources are uncertain.

% A practical barrier to two competitors sharing resources is that they 
% may not know the value for 
% each others' resources until recieving and using them.
% For example if two competitors are exchanging different datasets, a competitor may not know their 
% value for the other's data until receiving it and attempting to use it to train a machine learning model.
% In the case of research labs exchanging access to compute or equipment, the value for this access 
% could be contingent upon the results of experiments or simulations run on these resources. If an
% experiment/simulation produced novel results then access was valuable but if it did not, then the access 
% was substantially less valuable.
% The main question of this work is how can we incentivize competing agents to share resources
% when it is likely mutually beneficial to do so, given uncertainty about their valuations for each
% others' resources.

% goal is to fin

% \begin{enumerate}

%     \item 
%     Give data examples

%     \item 
%     You don't know the value of compute time until you try to use it to do your experiments or 
%     simulations. If the simulations or experiments finish early, great, if not then need more compute.
%     Or find a negative vs positive result from running scientific experiments or simulations. 
%     For example your access to a piece of equipment was supper valuable if you got great results
%     but not very valuable if you didn't.

%     \item 

% \end{enumerate}  

% \begin{enumerate}
%     \item Paragraph about rising importance of data, unique from econ perspective that
%     once collected it can be freely shared or replecated.
    
%     \item 
%     Data sharing platforms exist but make unrealistic assumptions, agents willing 
%     to give away data for free (in same company) because don't assume competition.
%     Also assume things like truthful reporting.

%     \item 
%     Prior work has studied incentivizes in data sharing (truthfulness, fairness, )

%     \item In this work we study competition for strategic agents
    
% \end{enumerate}  

\subsection{Problem Motivation}

We will suppose that there are two competitors, $\competitora$ and $\competitorb$, each possessing 
collections of $\numdatab$ and $\numdataa$ 
different goods
respectively. Suppose that both agents would derive value from each others' goods based on the quantity
received.
Let this be quantified by random value functions
$\vala:\cbr{0,\ldots,\numdataa}\rightarrow\RR$ and 
$\valb:\cbr{0,\ldots,\numdatab}\rightarrow\RR$, where $\vala$ and $\valb$ are independent.


\begin{restatable}{example}{nme-val-functions}
    \label{ex:nme-val-functions}
    Suppose $\competitora$ and $\competitorb$ possess datasets
    \(
        \cbr{X^\competitora_i}_{i=1}^\numdatab
        \drawniid
        % \Ncal\rbr{\mu_\competitora,\sigma^2}
        D_a
    \)
    and
    \(
        % \quad
        % \text{and}
        % \quad
        \cbr{X^\competitorb_i}_{i=1}^\numdatab
        \drawniid
        % \Ncal\rbr{\mu_\competitorb,\sigma^2}
        D_b
    \)
    drawn from bounded distributions
    % where $D_a$ and $D_b$ are bounded.
    % respectively
    but wish to estimate the mean of each others' distributions. In this case, the value 
    functions may be defined as
    \begin{align*}
    % \(
        \vala_{\numdatafora}
        =
        % \rbr{
        \bigg(
            \frac{1}{k}
            \sum_{k=1}^\numdatafora
            X^b_k
            -
            \mu_b
        \bigg)^2
        % ~\forall \numdatafora\geq 1
        \quad
        \text{and}
        \quad
        % }^2
    % \),
    % \(
        \valb_{\numdataforb}
        =
        % \rbr{
        \bigg(
            \frac{1}{k}
            \sum_{k=1}^\numdataforb
            X^\competitora_k
            -
            \mu_\competitora
        \bigg)^2
        % ~\forall \numdataforb\geq 1
        % }^2
    % \),
    \end{align*}
    for $\numdatafora,\numdataforb\geq 1$, and $\vala_0=-C_a$, $\valb_0=-C_b$ 
    for sufficiently large $C_a,C_b\in\RR$,
    representing agents' values when they lack data.
\end{restatable}
% Beyond Example~\ref{ex:nme-val-functions} 

More generally, in machine learning, the value for data may be unknown until using it to train a model and
observing its performance. 

However, the agents may be hesitant to exchange resources if they are in competition. 
A exchange 
\(
    \rbr{\exchangea,\exchangeb}
    =
    \exchange\in\setexchanges
    =
    \cbr{0,\ldots,n}\times\cbr{0,\ldots,m}
\)
corresponds to $\competitora$ receiving $\numdatafora$ of $\competitorb$'s resources and 
$\competitorb$ receiving $\numdataforb$ of $\competitora$'s resources.
For a particular exchange 
\(
    \rbr{\exchangea,\exchangeb}
    =
    \exchange\in\setexchanges
    % =
    % \cbr{0,\ldots,n}\times\cbr{0,\ldots,m}
\)
let
\(
    \utilas = \utila\big(\valai,\valbj\big)
\)
and
\(
    \utilbs = \utilb\big(\valai,\valbj\big)
\)
% let $\utila:\setexchanges\rightarrow\RR$ and $\utilb:\setexchanges\rightarrow\RR$ denote
be the utilities for both agents,
where
\(
    \utila,\utilb:\RR^2\rightarrow\RR
\)
are functions mapping the value obtained by both competitors to their utilites.
In other words, $\utila,\utilb$ model how competitors care about either value 
derived from collaborating.

For example, a natural model is when competitors' utilies are linear both part parties'
valuations and are given by
\(
    \utilas
    =
    \valai-\compcoeff\valbj
\)
and
\(
    \utilbs
    =
    \valbj-\compcoeff\valai
\)
for $\compcoeff\in(0,1)$.
In this model, $\compcoeff$ quantifies how directly the agents are in competiton.
When $\compcoeff=1$ agents are in a zero-sum game whereas $\compcoeff=0$
signifies no competition. 

Let $\welfarefunc:\RR^2\rightarrow\RR$ be a selected welfare function used to quantify the quality of an exchange
$\exchange\in\setexchanges$ via $\welfarefunc\rbr{\utilas,\utilbs}$. 
Common examples of welfare functions include the egalitarian welfare~\citep{mas1995microeconomic} 
$\welfarefunc(x,y)=\min(x,y)$, the 
Nash welfare~\citep{nash1950bargaining} $W(x,y)=xy$, or the social welfare subject to individual rationality 
$\welfarefunc(x,y)=(x+y)\cdot 1_{\cbr{x,y\geq 0}}$. We are interested in the problem of designing  
algorithm which facilitates an exchange $\exchange\in\setexchanges$ to maximize the welfare.

If the valuation functions $\vala,\valb$ are deterministic, this is not an interesting problem as it simply reduces to 
the discrete optimization problem
\(
    \argmax_{
        \exchange\in\setexchanges
    }
    \welfarefunc\rbr{
        \utilas,\utilbs
    }
\). 
When $\vala$ and $\valb$ are random,
% However, agents may not know their value for a particular quantity of goods until after 
% receiving and using them? 
% For example, an agent may value data based on the accuracy of a model trained from it, which is 
% unknown until the data is used. In such cases what exchanges should a data sharing algorithm propose 
% to maximize the welfare?
% To model this, we assume that $\vala$ and $\valb$ are independent random variables, each drawn from known but not 
% necessarily identical distributions. 
% Under this model 
a natural approach would be to propose an
exchange that maximizes the expected welfare, i.e. to solve for 
% \begin{align*}
\(
    \argmax_{
        \exchange\in\setexchanges
    }
    \EE\sbr{
        \welfarefunc\rbr{
            \utilas,\utilbs
        }
    } .
\)
% \end{align*}

However, it turns out that if instead of proposing a one time exchange, we propose a sequence of increasing
exchanges (with both competitors reporting their observed values after each round), then we can obtain a 
higher expected welfare. This is demonstrated in the following example.

\begin{restatable}{example}{multi-vs-single-round-alg}
    \label{ex:multi-vs-single-round-alg}
    Let $\numdataa=\numdatab=1$, so $\setexchanges=\cbr{0,1}^2$. Suppose that each competitor's value for the other's good is 
    $\valaii{1},\valbjj{1}\drawniid\Bern(1/2)$ and assume $\valaii{0}=\valbjj{0}=0$.
    Let the utilities and welfare be given by
    \begin{align*}
        \utilas
        =
        \valai-\frac{1}{4}\valbj,
        \quad 
        \utilbs
        =
        \valbj-\frac{1}{4}\valai,
        \quad
        \text{and}
        \quad
        \welfarefunc(x,y)=\min(x,y) .
    \end{align*}  
    A direct computation shows that 
    \begin{align*}
        \rbr{1,1} 
        =
        \argmax_{
            \exchange\in\cbr{0,1}^2
        }
        \EE\sbr{
            \min\rbr{
                \utilas,\utilbs
            }
        } 
        \quad
        \text{and}
        \quad
        \EE\sbr{
            \min\rbr{
                \utilass{(1,1)},\utilbss{(1,1)}
            }
        } 
        =
        \frac{1}{16} .
    \end{align*}  
    Thus $(1,1)\in\setexchanges$ is the single exchange that maximizes the expected welfare.

    Now consider a multi-exchange algorithm that works as follows. The algorithm first gives one of 
    $\competitora$ or $\competitorb$ the other's resouce (without loss of generality assume 
    $\competitora$ receives $\competitorb$'s resource). 
    Competitor $\competitora$ is then asked to report their value of the received resource
    (either $0$ or $1$). If the $\competitora$ reports a $1$, then $\competitorb$ is 
    given $\competitora$'s resource. 
    However, if $\competitora$ reports a $0$, then $\competitorb$ is given nothing.
    Let $\exchange'=(i',j')\in\cbr{0,1}^2$
    denote the final exchange in this algorithm. A direct computation shows that 
    \begin{align*}
        \EE\sbr{
            \min\rbr{
                \utilass{s'},\utilbss{s'}
            }
        } 
        &=
        \EE\sbr{
            \min\rbr{
                \valaii{i'}
                -
                \frac{1}{4}
                \valbjj{j'}
                ,
                \valbjj{j'}
                -
                \frac{1}{4}
                \valaii{i'}
            }
        } 
        \\
        &=
        \frac{1}{2}
        \rbr{
            \frac{1}{2}
            \min\rbr{
                1-\frac{1}{4}
                ,
                1-\frac{1}{4}
            }
            +
            \frac{1}{2}
            \min\rbr{
                1-\frac{1}{4}
                ,
                -\frac{1}{4}
            }
        }
        +
        \frac{1}{2}\min(0,0)
        \\
        &=
        \frac{1}{8} .
    \end{align*}  
\end{restatable}  

Example~\ref{ex:multi-vs-single-round-alg} demonstrates that algorithms which can use competitors'
reported valuations to propose an increasing sequence of exchanges can obtain a strictly higher
expected welfare than a single exchange algorithm. Thus our goal will be to find the
optimal multi-exchange algorithm.

\textbf{Truthfulness assumptions.}
% Multi-exchange algorithms like the one given in example 1.1 rely on competitors truthfully reporting their value
% or their value being publicly observable/validated 
% for the data they have received. This is because in order to fairly determine future exchanges in both
% competitors best interests, they need to know the utilities of both agents. 
% To capture this difficulty, we first define the price of stability in this problem to be the 
% ratio of the expected welfare of the optimal multi-exchange algorithm to the expected welfare of the 
% optimal truthful mult-exchange algorithm. We then present a hardness result in Prop(\mycomment{fill later}),
% showing that the price of stability is unbounded. In other words, truthful algorithms cannot compete 
% with multi-exchange algorithms.
Multi-exchange algorithms, like Example~\ref{ex:multi-vs-single-round-alg}, rely on competitors truthfully 
reporting (or publicly validating) the value of the resources they receive. This is essential for 
fairly determining future exchanges based on both agents' utilities. To formalize this challenge, 
we will define the price of stability as the ratio between the expected welfare of the optimal 
multi-exchange algorithm and that of the optimal truthful one. 
In Proposition (\mycomment{fill later}), we show this ratio is unbounded, i.e. truthful algorithms cannot 
match the performance of unrestricted multi-exchange algorithms.


\mycomment{say other person's utility is a function of their valuation???}

\mycomment{add an unbounded PoS result}
\mycomment{does PoS even make sense cause welfare can be 0 or negative?}
\mycomment{abs difference unbounded}
% \mycomment{put in here something saying in order to create multi round }

\textbf{When participants can exit early.}


\subsection{Results and Technical Contributions.}
Our main contributions are formulating the problem of incentivizing competing agents to exchange resources
given uncertainty about their valuations for each others' resources and the design of an optimal 
algorithm to do so.

\textbf{Our novel setting.} 
\mycomment{edit this to be a mini prior work of why our problem setting is different and not captured
by existing work}
\mycomment{argue that existing models in other papers don't capture increasingly important usecases like
data sharing for ml}
\mycomment{while doing this, introduce our notation}
In Section~\ref{sec:model} we start by formalizing the interaction protocol
between competitors when exchanging goods, defining their valuations, utilities, and the welfare under
a particular exchange. 
Motivated by Example~\ref{ex:multi-vs-single-round-alg}, we define an exchange algorithm, to be a sequence
of increasing proposed exchanges, each based on reported valuations from the prior exchanges.
From this, we define an optimal exchange algorithm to be one which maximizes the
expected welfare. 
Finding an algorithm is our first goal.
We also define what it means for an exchange algorithm to be individually rational: in expectation 
competitors are no worse off participating than forgoing collaboration.
\mycomment{make sure we introduce utility, welfare here to, cause mentioned in following section}

\textbf{Reducing algorithm design to optimal stopping.} 

\emph{Challenges:}
A central challenge in designing an optimal exchange algorithm is proposing a sequence 
of optimal exchanges using only the information revealed by the previous exchange about 
future exchanges.
This is further complicated by the lack of assumptions we require of 
the valuation functions,
utilty functions, and welfare function. For example, 
% for $\arbitrarycompetitor\in\cbr{\competitora,\competitorb}$, 
the random variables 
\(
    \big\{
        \valai
    \big\}_{\numdatafora=0}^\numdataa
\)
could be continuous, discrete, or arbitrarily correlated
as is also the case with
\(
    \big\{
    % \{
        \valbj
    \big\}_{\numdataforb=0}^\numdatab
    % \big\}_{\numdataforb=0}^\numdatab
\).
Furthermore, the only assumption we make on these random variables (and the utilty and welfare funcitons 
$\utilas,\utilbs,\welfarefunc$) is boundedness.
Because of this, projecting the welfare of future exchanges is complicated.
Thus, it is difficult to design an algorithm which repeatedly proposes
a sequence of exactly optimal exchanges given the available information from the previous exchanges.
Competition adds another layer of complication as exhanges that are in one competitor's best interest may not 
be in the interest of the other's.

\emph{Techniques and insights:}
Our first insight is to reframe an algorithm as a policy for traversing the lattice 
\(
    \setexchanges
    =
    \cbr{0,\ldots,\numdataa}
    \times
    \cbr{0,\ldots,\numdatab}
\), 
which encodes the set of all valid exchanges,
starting at $(0,0)$ and moving only up/right at each step, using the information 
reported from previous exchanges.
Each such traversal corresponds to an increasing sequence of exchanges, representing that competitors
cannot unshare resources.
Our core insight is that we can use this reformulation to reduce the design of an optimal exchange
algorithm to solving an optimal stopping time problem over the partially ordered lattice $\setexchanges$,
leveraging techniques from~\citep{mandelbaum1981optimal}.
At a high level, the algorithm works as follows. At each time, determine which of the following three options would lead to 
a higher expected welfare: (1) terminating the exchange, (2) giving $\competitora$ an additional one of 
$\competitorb$'s resources and then ``proceeding optimally'', or (3) giving $\competitorb$ an additonal one of $\competitorb$'s
resources and then ``proceeding optimally''.
We compute the expected welfare of ``proceeding optimally'' after a given exchanged, using a dynamic programming 
type result from~\citep{mandelbaum1981optimal} which is a common approach in optimal stopping.
We show that under a reasonable assumption on the welfare function $\welfarefunc$,
the optimal algorithm we obtain is also automatically individually rational.

% \begin{enumerate}
%     \item A main challenge in designing an opt alg is that the space of all possible algorithms is extremely complicated
%     and that we do not put \textbf{any} restrictions on the valuation functions, for example they could be discrete support
%     or they could be continues and markovian, give example of trending upwards markov normal dist model? We also do not 
%     make any assumptions on the welfare functions, utility funcitons, other than normalization so that when no data is shared
%     welfares and utilities are 0. Competiton also inhererntly makes this complicated because exchanges that are in one agents' 
%     best interest may not be in the others'. 
    
%     \item
%     Our first technical insight is to reframe the design of an exchange algorithm to the design of an algorithm which traverses
%     the lattice (S= ..., moving up/right) which encodes the possible exchanges between agents. Our core insight is that 
%     we can use this reformulation to reduce the design of an optimal exchange algorithm to solving an optimal stopping time 
%     over the partially ordered lattice using techniques from~(). 

%     \item
%     To perform this reduction we consider the state space ... along with the filtration ... to capture the information know
%     to the algorithm designer after both agents report their valuations after an exchange. 
%     The result is an algorithm which achieves the optimal expected revenue among all algorithms. At a high level, this algorithm
%     works as follows. At each time t give only one of either a or b the others' next good and ask them to report their valuation
%     based on this the algorithm can decide which agent is 
    
%     \mycomment{mention dynamic programming}
% \end{enumerate}  


% In Section~\ref{sec:opt-non-ic-alg} we show how the problem of 
% designing an optimal exchange algorithm can be reduced to solving an optimal stopping time problem over 
% the partially ordered lattice $\NN^2_{\geq 0}$ using the framework of~\citet{mandelbaum1981optimal}. 
% % stopping time problem over the partially ordered set $\NN^2_{\geq 0}$ (with $(x_1,y_1)\leq (x_2,y_2) \iff$ 
% % $x_1\leq x_2$ and $y_1\leq y_2$).
% % More concretely, we leverage an extension of the classical optimal stopping 
% Because an exchange algorithm is an increasing sequence of exchanges (we use the standard partial ordering
% over $\NN^2$ to compare exchanges: $(i_1,j_1)\leq(i_2,j_2)\iff i_1\leq i_2$ and $j_1\leq j_2$), solving 
% for the optimal exchange algorithm is equivalent to determinining how to traverse $\NN^2_{\geq 0}$ in an
% increasing manner and determining when to stop traversing (i.e. when to settle on a final exchange).
% In addition to proving optimality, we show that for any choice of welfare function that respects individual
% rationality, our optimal exchange algorithm is also individually rational.
% As our definition of individual rationality is ex-ante, we supplement this with high probability guarantees
% that our optimal exchange algorithm is ex-post individually rational.

\mycomment{sep line}

\mycomment{put in IC motivating example here? or is it weird to have an example in the contributions section}
\textbf{An optimal incentive compatible algorithm.}
Our model thus far has assumed that once an agent signs on to participate in an algorithm, it will 
adhere to the sequence of exchanges proposed until the algorithm terminates. 
However, a competitor may be motivated to exit the exchange process early if she gets lucky with the value
she obtains from early on exchanges. In Section~\ref{sec:opt-ic-alg} we define incentive compatible
exchange algorithms to be those where neither competitor ever benefits from exiting the 
exchange algorithm early. We then give simple examples where the optimal incentive compatible exchange algorithm 
achieves a strictly lower expected welfare than the optimal exchange algorithm, demonstrating that 
incentive compatiblility is a property that cannot be obtained for free. We will then give examples showing the 
different challenges at play that complicate the design of an optimal incentive compatible algorithm. 

A priori,
the answer to whether an optimal incentive compatible algorithm even exists is not clear. We answer this question
in the affirmative, by constructing such an algorithm, providing novel adaptations of concepts from the theory 
of optimal stopping. 
We show that in addition to being incentive compatible,
this algorithm satisfies the same ex-ante and ex-post individual rationality guarantees as our first algorithm.
% One caveat is that 


% \paragraph{Truthfulness impossibility results}


\subsection{Related Work}

\begin{enumerate}
    \item \mycomment{Is there existing econ literature on the exchange of goods under competition?}
    \item Look at peter's? for info
    \item related-work-planning.md
\end{enumerate}  